\documentclass[10pt,a4paper]{amsart}
\usepackage[margin=19mm]{geometry}
\usepackage{amsmath,amssymb,mathtools}
\usepackage[scr=rsfs,cal=euler]{mathalfa}
\usepackage{xfrac}
\usepackage[unicode,hidelinks,bookmarks=false]{hyperref}
\newcommand{\ie}{i.e.\ }
\newcommand{\cf}{cf.\ }
\newcommand{\resp}{respectively\ }
\newcommand{\Subsection}{\subsection}
\providecommand{\nobreakdash}{\nobreak\hbox{-}}
\setlength{\emergencystretch}{2em}
\multlinegap0pt
\usepackage{braket}
\usepackage{amssymb}
\usepackage{dsfont}
\usepackage{tikz}
\newtheorem{theorem}{Theorem}
\title{On the intrinsic geometry of polyhedra: Convex polygon coordinates}
\author{Anna B. Romanowska, Jonathan D.H. Smith, Anna Zamojska-Dzienio}
\date{Source: arXiv:2603.08757v1 (CC BY 4.0)}
\begin{document}
\maketitle

\noindent\emph{Source and licence.} On the intrinsic geometry of polyhedra: Convex polygon coordinates. Version arXiv:2603.08757v1 (CC BY 4.0), distributed by arXiv under the Creative Commons Attribution 4.0 International licence, http://creativecommons.org/licenses/by/4.0/, as recorded in the arXiv licence metadata for this exact version and shown on the abstract page https://arxiv.org/abs/2603.08757v1. Full licence notice: \texttt{licenses/arxiv-2603.08757-CC-BY-4.0.txt}.\par\smallskip

\noindent\textit{Editorial scope.} Counted unit: the article's theorem T:PolCoSys (source0.tex lines 510--512). The article's complete original proof of it is delivered with it (source0.tex lines 515--555, one \texttt{proof} environment of 41 lines). No mathematics is written, repaired or re-derived: the statement and the proof are the article's own lines, in the article's own notation, and no retained line is substituted or re-flowed.  Retained uncounted as the source-local context the proof needs: lines 469--471; lines 494--496.  Nothing was omitted from any retained span, so the package carries no omission record.  No retained line contains a citation, so the wrapper carries no bibliography and no bibliography entry.  No retained line contains a figure environment, an image-inclusion command or drawing code, so no image is delivered and the package declares no asset dependency.  Editorial formatting only, in addition: an AMS \texttt{amsart} preamble that restates the article's own declarations and macros used by the retained lines, namely the environments \texttt{theorem} on separate counters, together with the standard packages those lines need.  The retained environments print in this document as Theorem 1 in order of appearance, whereas the article prints the same items with the numbering of its own full document.  The complete native source of the article as distributed in the arXiv e-print for this exact version is bundled unchanged as \texttt{sources/196028/source0.tex}.
\begin{equation}\label{E:BraktPid}
\sum_{v\in V}\ket v\bra v=1_\Pi
\end{equation}

\begin{equation}\label{E:CoSyKapa}
\lambda\colon V\to I^\Pi;\mathbf v\mapsto\ket{\mathbf v}_\lambda
\end{equation}

\begin{theorem}\label{T:PolCoSys}
The set $K_\Pi$ of cooordinate systems on a polytope $\Pi$ with vertex set $V$ forms a convex subset of the hypercube $I^{\Pi\times V}$ under pointwise barycentric operations.
\end{theorem}

\begin{proof}
Using the notation of \eqref{E:CoSyKapa}, a coordinate system $\lambda$ is specified by the element $(\ket{\mathbf v_1}_\lambda,\dots,\ket{\mathbf v_n}_\lambda)$ of the pointwise barycentric algebra $(I^\Pi)^V$. Now, consider coordinate systems $\lambda,\lambda'$ and $p\in I^\circ$. By \eqref{E:BraktPid}, we have
\begin{equation}\label{E:lamblamb}
\sum_{\mathbf v\in V}\ket{\mathbf v}_{\lambda}\bra{\mathbf v}=1_\Pi
\quad
\mbox{and}
\quad
\sum_{\mathbf v\in V}\ket{\mathbf v}_{\lambda'}\bra{\mathbf v}=1_\Pi\,.
\end{equation}
The potential coordinate system $\lambda\lambda'\underline p$ is given by
$$
\ket{\mathbf v}_{\lambda\lambda'\underline p}
=\ket{\mathbf v}_{\lambda}\ket{\mathbf v}_{\lambda'}\underline p
$$
for $\mathbf{v}\in V$. Then
\begin{align*}
\sum_{\mathbf v\in V}\ket{\mathbf v}_{\lambda\lambda'\underline p}\bra{\mathbf v}
&
=
\sum_{\mathbf v\in V}\left(\ket{\mathbf v}_{\lambda}\ket{\mathbf v}_{\lambda'}\underline p
\right)\bra{\mathbf v}
=
\sum_{\mathbf v\in V}
\left(\ket{\mathbf v}_{\lambda}\bra{\mathbf v}\right)
\left(\ket{\mathbf v}_{\lambda'}\bra{\mathbf v}\right)
\underline p
\\
&
=
\left(
\sum_{\mathbf{v}\in V}\ket{\mathbf v}_{\lambda}\bra{\mathbf v}
\right)
\left(
\sum_{\mathbf{v}\in V}\ket{\mathbf v}_{\lambda'}\bra{\mathbf v}
\right)
\underline p
=1_\Pi1_\Pi\underline p
=1_\Pi
\end{align*}
using \eqref{E:lamblamb} for the penultimate equality. By Definition~\ref{T:PolCoSys}, it follows that $\lambda\lambda'\underline p$ is indeed a coordinate system.
\end{proof}

\end{document}
